% !TeX root = ../main.tex
\section{Threats to Validity}
\label{sec:threats-to-validity}

\subsection{Construct Validity (Partial)}
\label{sec:construct-validity}

\emph{Architectural information in requirements.}
User stories and acceptance criteria describe required behavior but may
omit quality attributes, deployment constraints or architectural context
needed to judge a pattern. System purpose supplies context, not sufficient
evidence of applicability by itself. The protocol retains an
insufficient-information category and requires requirement references for
positive classifications (Section~\ref{sec:applicability-experiment}).
These safeguards avoid forcing every pair into a binary decision, but
neither validate the adequacy of the supplied context nor explain the
dominant label distribution in Table~\ref{tab:applicability-results}.
Reference membership checks do not establish semantic entailment.
Independent review of the cited evidence is still needed; controlled
context enrichment is a future mitigation, not an intervention already
shown to improve classification.

\emph{Applicability, retrieval overlap and ADR quality.}
Applicability is operationalized as a model judgment under the supplied
prompt and catalog, not independently established architectural truth.
Final Top-5 overlap measures same-story candidate membership relative to
the LLM-positive set (Section~\ref{sec:rag-baseline}); it does not measure
accuracy, precision or recall. Explicit denominators and separate reporting
of classification, overlap and ADR completion limit conflation between
these constructs. ADR schema conformity and nonempty fields establish
structural completeness, not correctness, sufficient grounding or
comprehensive trade-offs (Section~\ref{sec:adr-results}).
The planned professional assessment addresses separate judgment dimensions,
but cannot convert generation completion into a validated quality measure.

\emph{Questionnaire coverage.}
The prepared Q1--Q5 items distinguish decision adequacy, grounding,
consequences, perceived record usefulness and pattern applicability
(Table~\ref{tab:expert-constructs}). This separation and the predeclared
item-specific analysis avoid treating different dimensions as an
unjustified composite score. Nevertheless, explanatory questionnaire
wording and preparation checks do not establish instrument validity or
consistent respondent interpretation. Perceived usefulness is not observed
practice improvement; the items do not directly measure implementation
feasibility, decision effort or deployment success.
Optional comments may contextualize judgments, but do not replace missing
measurements. Confirming the administered wording and case mapping,
examining actual coverage and retaining critical comments are planned
safeguards, not completed empirical validations.

\emph{Evaluation of selected pairs only.}
The prepared questionnaire presents LLM-positive pairs and their generated ADRs.
This permits assessment of those recommendations, not a catalog-wide
reference set. Rejected and insufficient-information pairs are outside
the prepared evaluation, so the design cannot establish missed applicable
patterns, false-negative rates or comparative selection accuracy.
Restricting claims to assessed cases is the current mitigation.
Independent assessment of non-positive pairs or all patterns for a bounded
subset is proposed future work (Section~\ref{sec:research-implications}),
not evidence supplied by the selected-pair questionnaire.

\noindent\textbf{TBD---administered-instrument and expert-evidence
reconciliation pending.} After Phase 3, verify actual wording, routing,
eligibility and coverage, document deviations from the prepared protocol,
and revise these construct limits against the supported findings.
No completed instrument validation or expert-based mitigation effect is
claimed in this partial draft.
